\chapter{Mathematical Induction}
\signature

\section{The well-ordering principle}

\begin{theorem}[Well-ordering principle, WOP]
Every nonempty subset of $\N$ has a least element.
\end{theorem}

This innocuous statement is the bedrock: it is logically equivalent to the
principle of induction, and powers ``minimal counterexample'' arguments ---
if a statement failed for some $n$, there would be a \emph{smallest} failing
$n$, which one then contradicts.

\section{Weak induction}

\begin{theorem}[Principle of mathematical induction, PWI]\label{thm:pwi}
Let $P(n)$ be a predicate on $\N$. If
\begin{itemize}
\item[(i)] \textbf{Base case:} $P(0)$ holds (or $P(n_0)$ for a starting
  point $n_0$), and
\item[(ii)] \textbf{Inductive step:} for all $n \geq n_0$,
  $P(n) \Rightarrow P(n+1)$,
\end{itemize}
then $P(n)$ holds for all $n \geq n_0$.
\end{theorem}

\begin{proof}[WOP implies PWI]
Suppose (i) and (ii) hold but $P$ fails somewhere. By WOP let $m$ be the
least failure. $m \neq n_0$ by (i), so $m - 1 \geq n_0$ and
$P(m-1)$ holds (minimality). By (ii), $P(m)$ holds --- contradiction.
(Conversely PWI implies WOP: see Exercise 13.10.)
\end{proof}

\textbf{The template.} To prove $\forall n \geq n_0,\ P(n)$:
(1) \emph{state} $P(n)$ clearly; (2) verify $P(n_0)$;
(3) fix $n \geq n_0$, assume the \emph{inductive hypothesis} $P(n)$, and
derive $P(n+1)$ --- pointing precisely at where the hypothesis is used;
(4) conclude by induction.

\section{Summation identities}

\begin{example}\label{ex:gauss}
$\displaystyle\sum_{i=1}^{n} i = \frac{n(n+1)}{2}$ for $n \geq 1$.
\emph{Base:} $1 = \frac{1 \cdot 2}{2}$. \emph{Step:} assuming the identity
for $n$,
\[
\sum_{i=1}^{n+1} i = \frac{n(n+1)}{2} + (n+1)
= (n+1)\Bigl(\frac{n}{2} + 1\Bigr) = \frac{(n+1)(n+2)}{2}. \qed
\]
\end{example}

\begin{example}
$\displaystyle\sum_{k=0}^{n} 2^{k} = 2^{n+1} - 1$:\ base $1 = 2 - 1$; step:
add $2^{n+1}$ to both sides: $2^{n+1} - 1 + 2^{n+1} = 2^{n+2} - 1$.
Similarly $\sum_{k=1}^{n} k \cdot k! = (n+1)! - 1$: step uses
$(n+1)!\,(n+2) - 1 - \ldots$ --- see Exercise 13.2. The sums
$\sum k^2 = \frac{n(n+1)(2n+1)}{6}$ and
$\sum k^3 = \bigl(\frac{n(n+1)}{2}\bigr)^2$ follow the identical pattern
(Chapter~9).
\end{example}

\section{Inequalities}

\begin{example}\label{ex:2nn2}
$2^{n} > n^{2}$ for all $n \geq 5$.
\emph{Base:} $2^5 = 32 > 25$. \emph{Step:} for $n \geq 5$, assume
$2^n > n^2$. Then
\[
2^{n+1} = 2 \cdot 2^{n} > 2n^{2} = n^{2} + n \cdot n
\geq n^{2} + 5n > n^{2} + 2n + 1 = (n+1)^{2},
\]
using $n \geq 5 \Rightarrow 5n > 2n + 1$. Note the base case $5$ is
essential: the claim is false for $n = 2, 3, 4$.
\end{example}

\begin{example}
$n! \geq 2^{n}$ for $n \geq 4$: base $24 \geq 16$; step:
$(n+1)! = (n+1) \cdot n! \geq (n+1)2^{n} \geq 2 \cdot 2^{n} = 2^{n+1}$
since $n + 1 \geq 2$. Two classical companions provable the same way:
Bernoulli's inequality $(1+x)^{n} \geq 1 + nx$ for $x \geq -1$, and
$\ln(n+1) \leq H_n = \sum_{i=1}^n \frac{1}{i}$.
\end{example}

\section{Divisibility}

\begin{example}\label{ex:div}
$3 \mid (n^{3} + 2n)$ for all $n \geq 0$.
\emph{Base:} $3 \mid 0$. \emph{Step:}
\[
(n+1)^{3} + 2(n+1) = n^{3} + 3n^{2} + 3n + 1 + 2n + 2
= (n^{3} + 2n) + 3(n^{2} + n + 1),
\]
a sum of two multiples of $3$ (the first by hypothesis).
\end{example}

\section{Sets and geometry}

\begin{example}[Subsets]
A set with $n$ elements has $2^{n}$ subsets. \emph{Base:}
$\powerset(\varnothing) = \{\varnothing\}$, of size $1 = 2^0$.
\emph{Step:} let $|A| = n + 1$, pick $x \in A$, $A' = A \setminus \{x\}$.
Subsets of $A$ split into subsets of $A'$ and those plus $x$:
$2^{n} + 2^{n} = 2^{n+1}$ by hypothesis.
\end{example}

\begin{example}[Triangulation]
Every convex polygon with $n \geq 3$ vertices can be divided into $n - 2$
triangles by non-crossing diagonals. \emph{Base:} a triangle, $1 = 3 - 2$.
\emph{Step:} in a convex $(n+1)$-gon, the diagonal from $v_1$ to $v_n$ cuts
off triangle $v_1 v_n v_{n+1}$ and leaves a convex $n$-gon, triangulated by
hypothesis into $n-2$ triangles: total $n - 1 = (n+1) - 2$.
\end{example}

\section{Strong induction}

\begin{theorem}[Strong induction]\label{thm:strong}
If $P(n_0)$ holds and, for all $n \geq n_0$,
\[
\bigl(P(n_0) \wedge P(n_0 + 1) \wedge \cdots \wedge P(n)\bigr)
\;\Rightarrow\; P(n+1),
\]
then $P(n)$ holds for all $n \geq n_0$. Strong and weak induction are
equivalent in power, but strong induction is the natural tool when $P(n+1)$
depends on values \emph{further back} than $n$.
\end{theorem}

\begin{example}[Prime factorisation]\label{ex:primefact}
Every integer $n \geq 2$ is a product of primes. If $n + 1$ is prime, done.
Otherwise $n + 1 = ab$ with $2 \le a, b \le n$; by the strong hypothesis
both $a$ and $b$ are products of primes, hence so is $n+1$.
(Weak induction is helpless here: $a, b$ are far below $n + 1$.)
\end{example}

\begin{example}[Fibonacci bound]
$F_n < 2^{n}$ for all $n \geq 0$: bases $F_0 = 0 < 1$, $F_1 = 1 < 2$; step
($n \ge 1$): $F_{n+1} = F_n + F_{n-1} < 2^{n} + 2^{n-1} < 2^{n+1}$, using
the hypothesis at \emph{two} previous points.
\end{example}

\begin{example}[Postage stamps]
Every amount $n \geq 12$ can be paid with $4$- and $5$-dirham stamps.
Bases: $12 = 4+4+4$, $13 = 4+4+5$, $14 = 4+5+5$, $15 = 5+5+5$. Step for
$n + 1 \geq 16$: by strong hypothesis $(n+1) - 4 \geq 12$ is payable; add
one $4$-stamp.
\end{example}

\section{Structural induction}

For recursively defined structures (Chapter~14) --- formulas, strings, trees
--- induction follows the structure: prove the property for the base
objects, and show each construction rule preserves it. Example: every wff of
propositional logic has as many left as right parentheses --- true for
atoms (none), and each formation rule adds one of each.

\section{Pitfalls}

\begin{example}[All horses are the same colour]
``Proof'' that any $n$ horses share a colour: for the step, remove one horse
--- the remaining $n$ share a colour by hypothesis; remove a different one
--- same conclusion; the overlap forces all $n + 1$ to match. The argument
\emph{fails at $n + 1 = 2$}: the two groups of one horse do not overlap, so
nothing links them. One broken link ($P(1) \to P(2)$) destroys the whole
chain.
\end{example}

\begin{notebox}
Always check: (1) the base case is actually verified (a missing or false
base ``proves'' $\sum_{i=1}^n i = \frac{n(n+1)}{2} + 1$ too, since the step
goes through!); (2) the inductive step is valid \emph{for every} $n$ from
the base upward, including the very first; (3) the inductive hypothesis is
genuinely used --- otherwise the proof is direct, or wrong.
\end{notebox}

\section{Induction and programs}

Induction is the mathematics of loops and recursion: proving that a
recursive program is correct \emph{is} an induction on the input (see
Chapter~14), and a loop is proved correct via a \emph{loop invariant} ---
a property maintained by every iteration, i.e.\ an inductive step. Example:
in the iterative factorial, the invariant ``after $k$ iterations the
accumulator equals $k!$'' is established at $k = 0$ and preserved each pass,
so on exit (at $k = n$) the result is $n!$.

\section{Summary}

WOP, weak induction and strong induction are equivalent foundations. The
template --- state $P(n)$, verify the base, derive $P(n+1)$ using the
hypothesis visibly --- proves summation identities, inequalities (mind the
true base case), divisibility facts, set-theoretic and geometric claims.
Strong induction handles recurrences reaching far back (prime
factorisation, Fibonacci); structural induction handles recursively built
objects; and loop invariants are induction wearing a programmer's coat.

\section*{Exercises}
\addcontentsline{toc}{section}{Exercises}

\begin{exercise}{\easy}
Prove by induction: $\sum_{i=1}^{n} (2i - 1) = n^{2}$ for $n \geq 1$.
\end{exercise}

\begin{exercise}{\easy}
Prove by induction: $\sum_{k=1}^{n} k \cdot k! = (n+1)! - 1$.
\end{exercise}

\begin{exercise}{\easy}
Prove by induction: $\sum_{k=1}^{n} \frac{1}{k(k+1)} = \frac{n}{n+1}$.
\end{exercise}

\begin{exercise}{\medium}
Prove that $n! > 3^{n}$ for all $n \geq 7$.
\end{exercise}

\begin{exercise}{\medium}
Prove that $6 \mid (n^{3} - n)$ for all $n \geq 0$.
\end{exercise}

\begin{exercise}{\medium}
Prove by induction that $\sum_{k=1}^{n} k^3 =
\bigl(\frac{n(n+1)}{2}\bigr)^2$.
\end{exercise}

\begin{exercise}{\medium}
Prove Bernoulli's inequality: $(1 + x)^{n} \geq 1 + nx$ for all real
$x \geq -1$ and $n \in \N$. Where does the step use $x \geq -1$?
\end{exercise}

\begin{exercise}{\medium}
A chocolate bar is an $m \times n$ grid of squares. Prove (strong induction
on the number of squares) that exactly $mn - 1$ breaks along grid lines are
needed to reduce it to unit squares.
\end{exercise}

\begin{exercise}{\hard}
Prove by strong induction that every $n \geq 1$ has a binary representation:
it can be written as $\sum_{i} b_i 2^{i}$ with $b_i \in \{0,1\}$.
\end{exercise}

\begin{exercise}{\hard}
Deduce the well-ordering principle from weak induction: if
$S \subseteq \N$ has no least element, prove by strong induction that
$S = \varnothing$.
\end{exercise}

\begin{exercise}{\hard}
Find the flaw: ``Theorem: $a^{n} = 1$ for every $a \neq 0$ and $n \geq 0$.
Base: $a^0 = 1$. Strong step:
$a^{n+1} = \frac{a^{n} \cdot a^{n}}{a^{n-1}} = \frac{1 \cdot 1}{1} = 1$.''
\end{exercise}

\section*{Solutions to the Exercises}
\addcontentsline{toc}{section}{Solutions to the Exercises}

\begin{solution}{13.1}
Base: $1 = 1^2$. Step: $\sum_{i=1}^{n+1}(2i-1) = n^2 + (2n+1) = (n+1)^2$.
\end{solution}

\begin{solution}{13.2}
Base: $1 \cdot 1! = 1 = 2! - 1$. Step:
$\sum_{k=1}^{n+1} k\,k! = (n+1)! - 1 + (n+1)(n+1)!
= (n+1)!\,(n+2) - 1 = (n+2)! - 1$.
\end{solution}

\begin{solution}{13.3}
Base: $\frac{1}{2} = \frac{1}{2}$. Step:
$\frac{n}{n+1} + \frac{1}{(n+1)(n+2)}
= \frac{n(n+2) + 1}{(n+1)(n+2)}
= \frac{(n+1)^2}{(n+1)(n+2)} = \frac{n+1}{n+2}$.
\end{solution}

\begin{solution}{13.4}
Base: $7! = 5040 > 2187 = 3^7$. Step: for $n \geq 7$,
$(n+1)! = (n+1)\,n! > (n+1)\,3^{n} \geq 8 \cdot 3^{n} > 3^{n+1}$.
\end{solution}

\begin{solution}{13.5}
Base: $6 \mid 0$. Step:
$(n+1)^3 - (n+1) = (n^3 - n) + 3(n^2 + n) = (n^3 - n) + 3n(n+1)$.
The first term is divisible by $6$ by hypothesis; $n(n+1)$ is even, so
$3n(n+1)$ is divisible by $6$ as well.
\end{solution}

\begin{solution}{13.6}
Base: $1 = 1$. Step: with $T_n = \frac{n(n+1)}{2}$,
$\sum_{k\le n+1} k^3 = T_n^2 + (n+1)^3
= \frac{(n+1)^2}{4}\bigl(n^2 + 4(n+1)\bigr)
= \frac{(n+1)^2 (n+2)^2}{4} = T_{n+1}^2$.
\end{solution}

\begin{solution}{13.7}
Base: $(1+x)^0 = 1 \geq 1$. Step:
$(1+x)^{n+1} = (1+x)^{n}(1+x) \geq (1 + nx)(1 + x)$ --- this multiplication
preserves the inequality \emph{because} $1 + x \geq 0$, i.e.\ $x \geq -1$.
Expanding: $1 + (n+1)x + nx^{2} \geq 1 + (n+1)x$.
\end{solution}

\begin{solution}{13.8}
$P(s)$: any bar of $s$ squares needs $s - 1$ breaks. Base $s = 1$: $0$
breaks. Step: a bar of $s + 1$ squares is first broken into pieces of $a$
and $b$ squares, $a + b = s + 1$, $a, b \le s$. By strong hypothesis they
need $a - 1$ and $b - 1$ further breaks: total
$1 + (a-1) + (b-1) = s = (s+1) - 1$, independent of the chosen break.
\end{solution}

\begin{solution}{13.9}
Base: $1 = 2^0$. Step: let $n + 1 \geq 2$. If $n + 1$ is even, write
$n + 1 = 2m$ with $1 \le m \le n$; by strong hypothesis
$m = \sum b_i 2^{i}$, so $n + 1 = \sum b_i 2^{i+1}$. If odd,
$n + 1 = 2m + 1$ with $0 \le m \le n$, and
$n + 1 = 1 + \sum b_i 2^{i+1}$, the new bit $b_0 = 1$.
\end{solution}

\begin{solution}{13.10}
Let $S \subseteq \N$ have no least element; set $P(n)$: ``no element of $S$
is $\leq n$''. $P(0)$: if $0 \in S$ it would be least, so $0 \notin S$.
Step: assume $P(n)$, i.e.\ $S \cap \{0,\dots,n\} = \varnothing$. If
$n + 1 \in S$ it would be the least element --- impossible; so $P(n+1)$.
By induction every $n$ avoids $S$: $S = \varnothing$. Contrapositive: every
nonempty $S$ has a least element.
\end{solution}

\begin{solution}{13.11}
The step at $n = 0$ computes $a^{1} = \frac{a^0 a^0}{a^{-1}}$, invoking the
hypothesis at $n - 1 = -1$, which is outside the induction's range (and the
base only covers $n = 0$). The chain $P(0) \Rightarrow P(1)$ is broken, so
nothing beyond the base is established.
\end{solution}
